\documentclass[12pt,a4paper]{article}

\usepackage{fontawesome5}
\providecommand{\faPencilAlt}{\faIcon{pencil-alt}}

\usepackage{styles/couleurs}
\usepackage{styles/mise_en_page}
\usepackage{styles/boites}
\usepackage{styles/maths}
\usepackage{styles/figures}

\usepackage{enumitem}
\usepackage{array}
\usepackage{tabularx}
\usepackage{amsmath}
\usepackage{amssymb}
\usepackage{tikz}
\usepackage{xparse}
\usepackage{graphicx}
\usepackage{pdflscape}
\usepackage{colortbl}

\usetikzlibrary{calc,arrows.meta,positioning,patterns}

\setlist[itemize]{leftmargin=1.4em,itemsep=0.15em,topsep=0.2em}
\setlist[enumerate]{leftmargin=*,itemsep=0.22em,topsep=0.25em}

\newcommand{\lignepointillee}{%
  \noindent\dotfill\par
  \vspace{0.45em}
}

\newcommand{\lignesreponse}[1]{%
  \begin{tcolorbox}[
    enhanced,
    colback=white,
    colframe=NeutreC!40,
    boxrule=0.5pt,
    arc=1.5mm,
    left=3mm,
    right=3mm,
    top=2mm,
    bottom=1mm,
    before skip=2pt,
    after skip=2pt
  ]
    \foreach \i in {1,...,#1}{\lignepointillee}
  \end{tcolorbox}
}

\newcommand{\imageactivite}[3]{%
  \IfFileExists{#3}{%
    \includegraphics[width=#1,height=#2,keepaspectratio]{#3}%
  }{%
    \fbox{%
      \parbox[c][#2][c]{#1}{%
        \centering\sffamily\small
        Image à intégrer\par
        \vspace{0.3em}
        \texttt{\detokenize{#3}}
      }%
    }%
  }%
}

\newcommand{\titreseq}[1]{%
\begin{center}
  {\sffamily\LARGE\bfseries\color{TitreC}
  Séquence 1 -- Statistiques}\\[0.2em]
  {\sffamily\large\bfseries\color{SavoirC}#1}
\end{center}
}

\begin{document}

\titreseq{Progression professeur}

\remarque{Durée indicative : sept séances d'apprentissage de 50 à 55 minutes,
deux interrogations courtes et une évaluation finale.

Les acquis de 5e ne sont pas repris sous forme d'une « liste de rappels » :
ils sont mobilisés comme des outils nécessaires pour résoudre des questions
courtes dès les deux premières activités.}

\section*{Place dans le programme}

En 4e, on poursuit le travail sur les effectifs, les fréquences,
la moyenne, les tableaux et les représentations graphiques.
Les nouveautés centrales de cette séquence sont :

\begin{itemize}
  \item la \textbf{médiane}, pour des séries de petit effectif total,
        pair ou impair ;
  \item la lecture, l'interprétation et la construction de
        \textbf{diagrammes circulaires}.
\end{itemize}

L'étendue et les histogrammes ne sont pas étudiés dans cette séquence :
ils relèvent de la suite de la progression du cycle.

\section*{Objectifs de fin de séquence}

L'élève doit pouvoir :

\begin{itemize}
  \item identifier population, caractère, valeurs ou modalités ;
  \item déterminer un effectif total et les effectifs de chaque valeur ;
  \item calculer une fréquence en écriture décimale ou en pourcentage ;
  \item organiser des données brutes dans un tableau ;
  \item choisir et lire une représentation adaptée ;
  \item construire un diagramme circulaire ;
  \item calculer et interpréter une moyenne simple ou à partir d'un tableau d'effectifs ;
  \item déterminer et interpréter une médiane pour un effectif pair ou impair ;
  \item comparer moyenne et médiane dans une situation comportant une valeur extrême ;
  \item rédiger quelques phrases pour communiquer ce que montrent les données ;
  \item contrôler un résultat ou une représentation.
\end{itemize}

\section*{Déroulé réaliste}

\renewcommand{\arraystretch}{1.38}
\begin{tabularx}{\textwidth}{
|>{\centering\arraybackslash}p{1.05cm}
|>{\raggedright\arraybackslash}p{4.2cm}
|>{\raggedright\arraybackslash}X
|>{\raggedright\arraybackslash}p{3.15cm}|}
\hline
\textbf{Séance} & \textbf{Support principal} &
\textbf{Enjeu didactique} & \textbf{Production}\\
\hline

1 &
Activité 1 : Quelle équipe choisir ? &
Faire comparer deux séries avant de nommer les outils.
Réactiver la moyenne et faire apparaître son insuffisance
lorsque deux équipes ont la même moyenne. Introduire la médiane. &
Argument individuel, débat en groupe puis trace courte.\\
\hline

2 &
Activité 2 : Mini data challenge &
Réactiver vocabulaire, effectifs, fréquences, tableaux et choix
de représentations à partir de données déjà recueillies.
Éviter le coût d'une enquête complète. &
Une mission par groupe puis mise en commun.\\
\hline

3 &
Activité 3 : Du tableau au disque &
Construire un diagramme circulaire en reliant effectifs,
fréquences, pourcentages et angles. &
Diagramme construit et contrôlé.\\
\hline

4 &
Activité 4 : Une moyenne sans recopier 29 valeurs &
Donner du sens au calcul d'une moyenne à partir d'un tableau
d'effectifs. &
Interrogation 1 en début de séance, puis calcul organisé.\\
\hline

5 &
Activité 5 : Moyenne ou médiane ? &
Comparer les deux indicateurs, étudier l'effet d'une valeur extrême
et choisir un résumé adapté au contexte. &
Justification écrite courte.\\
\hline

6 &
Activité 6 : Diagrammes et médias &
Développer des outils de contrôle : somme des pourcentages,
cohérence entre secteurs et nombres, titre, source, effectif. &
Interrogation 2 en début de séance, puis critique argumentée.\\
\hline

7 &
Consolidation &
Médiane à partir de données brutes, d'un tableau ou d'un diagramme
en bâtons ; fréquence, moyenne pondérée, interprétation. &
Automatismes et exercices différenciés.\\
\hline

8 &
Évaluation finale &
Mobiliser les outils dans des formes variées et expliquer les résultats. &
Évaluation sur 20.\\
\hline
\end{tabularx}

\section*{Organisation des deux premières activités}

\subsection*{Activité 1}

\begin{enumerate}
  \item Recherche individuelle : chacun prend position avec les outils de son choix.
  \item Mise en commun par groupes de trois : les élèves confrontent leurs critères.
  \item Débat bref : on constate que la moyenne seule ne permet pas de départager.
  \item Institutionnalisation : construction de la médiane et interprétation.
\end{enumerate}

\subsection*{Activité 2}

Les données sont déjà collectées et anonymisées.
Chaque groupe reçoit une mission différente, ce qui réduit le temps
de traitement tout en permettant une mise en commun riche :

\begin{itemize}
  \item organiser les réponses qualitatives ;
  \item organiser les réponses quantitatives ;
  \item calculer et interpréter des indicateurs ;
  \item rédiger un portrait statistique de la classe fictive.
\end{itemize}

\section*{Automatismes}

Rituel de cinq minutes, quatre à six questions :

\begin{itemize}
  \item séances 1 à 3 : effectif total, effectif d'une valeur,
        fréquence simple et lecture de tableau ;
  \item séances 2 à 5 : calcul et interprétation de moyenne ;
  \item séances 3 à 7 : classement d'une série et recherche de médiane ;
  \item séances 4 à 7 : moyenne à partir d'un tableau d'effectifs ;
  \item séances 3 à 7 : conversions entre pourcentage et angle ;
  \item après la séquence : une question statistique courte toutes les deux semaines.
\end{itemize}

\section*{Principes didactiques}

\begin{itemize}
  \item Les élèves rencontrent d'abord une question à résoudre,
        pas une définition à recopier.
  \item Plusieurs procédures sont autorisées lors des recherches.
  \item La rédaction intervient après la recherche :
        on améliore une production pour la rendre compréhensible.
  \item Les automatismes sont reliés à un obstacle réellement rencontré.
  \item Une représentation est choisie pour répondre à une question,
        pas seulement pour « faire un graphique ».
  \item Les élèves construisent des outils de contrôle :
        total des effectifs, somme des fréquences, somme des angles,
        ordre de grandeur de la moyenne, position de la médiane.
  \item Les valeurs extrêmes sont utilisées pour distinguer
        le rôle de la moyenne et celui de la médiane.
\end{itemize}

\end{document}
